\section{Two inputs and the bounded-map obstruction}\label{sec:reduction}

The construction separates a problem about functions on the ball from a
curvature estimate.  The first input is a strictly plurisubharmonic
function with controlled derivatives and an obstruction to a prescribed
holomorphic real-part bound.  The second turns its complex Hessian bounds
into a complete metric with uniformly negative real sectional curvature.
We state both inputs precisely and then prove their implication for
Theorem~\ref{main:theorem}.

\subsection{Centered charts and the two propositions}

Let $\B=\{z\in\C^2:|z|<1\}$ and put
\[
 \psi(z)=-\log(1-|z|^2),\qquad
 I=(\psi_{i\bar j}).
\]
Here a complex Hessian is also viewed as a Hermitian form, linear in its
first argument.  To center coordinates at a point $z_0\in\B$, let
$r=|z_0|$, choose a unitary map $U$ with $U(r,0)=z_0$, and set
\begin{equation}\label{base:charts}
 T=U\circ T_r,\qquad
 T_r(\zeta)=\left(
 \frac{r+\zeta_1}{1+r\zeta_1},
 \frac{\sqrt{1-r^2}\,\zeta_2}{1+r\zeta_1}
 \right).
\end{equation}
These maps are ball automorphisms with $T(0)=z_0$.  The identity
\[
 1-|T_r(\zeta)|^2
 =\frac{(1-r^2)(1-|\zeta|^2)}{|1+r\zeta_1|^2}
\]
shows that $T^*I$ has the same expression in $\zeta$ as $I$ has in $z$:
the additional terms in $\psi\circ T$ are constant or pluriharmonic.
In particular, this matrix is the identity at $0$ and has zero first
derivatives there.  We call \eqref{base:charts} the centered ball charts.

For a smooth real function $\phi$ on $\B$, write
$h=\partial\bar\partial\phi$ and, in a centered chart $T$, write
\[
 h^T_{i\bar j}(\zeta)
 =\partial_{\zeta_i}\partial_{\bar\zeta_j}(\phi\circ T)(\zeta).
\]
The uniform control needed below is the existence of constants
$0<c\le C<\infty$ such that, for every centered chart,
\begin{equation}\label{base:control}
 c\,\mathrm{Id}\le h^T(0)\le C\,\mathrm{Id}
\end{equation}
and
\begin{equation}\label{base:jets}
 \left|\partial_k h^T_{i\bar j}(0)\right|
 +\left|\partial_{\bar\ell}h^T_{i\bar j}(0)\right|
 +\left|\partial_k\partial_{\bar\ell}h^T_{i\bar j}(0)\right|
 \le C
 \qquad(i,j,k,\ell\in\{1,2\}).
\end{equation}
The derivatives here are ordinary derivatives of the displayed component
functions.  In invariant form, \eqref{base:control} says $cI\le h\le CI$.

\begin{proposition}[A controlled potential]\label{base:existence}
There exist $\phi\in C^\infty(\B,\R)$ and a finite constant $C$ for
which \eqref{base:control} and \eqref{base:jets} hold with $c=1$, and
such that no $H\in\mathcal O(\B)$ and $C_H\in\R$ satisfy
\begin{equation}\label{base:no-minorant}
 \operatorname{Re}H(z)
 \le C_H+4\log\frac1{1-|z|}+\phi(z)
 \qquad(z\in\B).
\end{equation}
\end{proposition}

For any smooth real $\phi$ on $\B$, define its Hartogs domain by
\begin{equation}\label{metric:domain}
 M_\phi=\{(z,w)\in\B\times\C:t(z,w)<1\},
 \qquad t(z,w)=e^{\phi(z)}|w|^2.
\end{equation}
For $\lambda>0$ define the potential and its complex Hessian by
\begin{equation}\label{metric:potential}
 \mathcal K_\lambda=\lambda\psi-\log(1-t),\qquad
 \widehat g_\lambda=\partial\bar\partial\mathcal K_\lambda.
\end{equation}
Our real metric convention is
$g_\lambda=2\operatorname{Re}\widehat g_\lambda$: if a real vector $X$
has $(1,0)$ part $u$, then
$g_\lambda(X,X)=2\widehat g_\lambda(u,u)$.

\begin{proposition}[Hartogs metric transfer]\label{metric:transfer}
Let $\phi\in C^\infty(\B,\R)$ satisfy \eqref{base:control} and
\eqref{base:jets} for fixed $0<c\le C<\infty$.  There is a finite
$\lambda_0=\lambda_0(c,C)>0$ such that, for every
$\lambda\ge\lambda_0$, the domain $M_\phi$ in
\eqref{metric:domain} is contractible and the tensor
$g_\lambda=2\operatorname{Re}\widehat g_\lambda$ defined by
\eqref{metric:potential} is a smooth geodesically complete K\"ahler
metric.  For this fixed $\lambda$ there are constants
$0<A_\lambda\le B_\lambda<\infty$ such that
\[
 -B_\lambda\le K_{g_\lambda}(\sigma)\le-A_\lambda<0
\]
for every point of $M_\phi$ and every real two-plane $\sigma$ in its
tangent space.
\end{proposition}

The proof of Proposition~\ref{base:existence} occupies the analytic
construction below.  The global geometry and curvature calculation in
Sections~\ref{sec:hartogs}--\ref{sec:pinching} prove
Proposition~\ref{metric:transfer}.  We first establish the short argument
that connects their conclusions to bounded holomorphic maps.

\subsection{The Jacobian on the zero section}

\begin{lemma}\label{map:obstruction}
Let $\phi\in C^\infty(\B,\R)$ have the nonexistence property in
Proposition~\ref{base:existence}: no holomorphic $H$ and real constant
$C_H$ satisfy \eqref{base:no-minorant}.  Then no bounded holomorphic
map $F:M_\phi\to\C^3$ has a nowhere-vanishing complex Jacobian
determinant.
\end{lemma}

\begin{proof}
Suppose that $F=(F_1,F_2,F_3)$ is such a map, and choose $B\ge1$ so
that $|F_\ell|\le B$ on $M_\phi$ for every component.  Since
$M_\phi$ is open in $\C^3$, its coordinates $(z_1,z_2,w)$ are global.
The function
\[
 d(z)=\det\frac{\partial(F_1,F_2,F_3)}{\partial(z_1,z_2,w)}(z,0)
\]
is holomorphic and nowhere zero on $\B$.

Fix $z\in\B$ and put $\delta=1-|z|$.  The restriction
$F_\ell(\,\cdot\,,0)$ is bounded and holomorphic on the entire base
ball.  For $j=1,2$, the disk
$\{z+\zeta e_j:|\zeta|<\delta\}$ is contained in that ball.  Cauchy's
estimate, applied on disks with radii tending to $\delta$, gives
\begin{equation}\label{map:base-cauchy}
 |\partial_{z_j}F_\ell(z,0)|\le B\delta^{-1}.
\end{equation}
For this fixed $z$, the fiber in $M_\phi$ is the disk of radius
$e^{-\phi(z)/2}$.  Cauchy's estimate on this disk gives
\begin{equation}\label{map:fiber-cauchy}
 |\partial_wF_\ell(z,0)|\le B e^{\phi(z)/2}.
\end{equation}
In particular, the base estimate requires no uniform fiber radius near
$z$; it is taken entirely on the zero section.

The six terms in the determinant expansion, together with
\eqref{map:base-cauchy} and \eqref{map:fiber-cauchy}, yield
\[
 |d(z)|\le6B^3\delta^{-2}e^{\phi(z)/2},
\]
and hence
\begin{equation}\label{map:determinant-bound}
 \log|d(z)|^2
 \le 2\log6+6\log B+4\log\frac1{1-|z|}+\phi(z).
\end{equation}
Because $\B$ is simply connected and $d$ has no zeros, it has a global
holomorphic logarithm $\ell$.  For example, a primitive of the closed
holomorphic one-form $d^{-1}\partial d$, with a suitable constant,
satisfies $e^\ell=d$.  The holomorphic function $H=2\ell$ has
$\operatorname{Re}H=\log|d|^2$, so \eqref{map:determinant-bound}
contradicts \eqref{base:no-minorant}.
\end{proof}

\begin{proof}[Proof of Theorem~\ref{main:theorem}]
Choose $\phi$ from Proposition~\ref{base:existence}, and then choose
$\lambda$ as in Proposition~\ref{metric:transfer}.  The domain
$M=M_\phi\subset\C^3$ is contractible and carries the asserted smooth
complete K\"ahler metric with two uniform negative real-sectional
bounds.  Lemma~\ref{map:obstruction} excludes a bounded holomorphic
map to $\C^3$ with nowhere-vanishing Jacobian.  A biholomorphism to a
bounded domain would be such a map: the chain rule applied to its
holomorphic inverse makes its Jacobian invertible everywhere.
\end{proof}

The logarithm argument uses nonvanishing along the entire zero section.
It gives no assertion that a bounded map's determinant vanishes
identically.  Nor are all bounded holomorphic functions constant:
$z_1$ and $z_2$ are nonconstant bounded holomorphic functions on $M_\phi$.
